\documentclass{article}
\usepackage[utf8]{inputenc}
\usepackage{amsmath, amssymb, amsthm, mismath, ulem, gensymb, mathtools, nccmath, hyperref, tikz, pgfplots}
\usepackage[miktex]{gnuplottex}
\usepackage{lipsum}
\usepackage{xcolor}
\usepackage{etoolbox}
\usepackage{hyperref}
\usepgfplotslibrary{fillbetween}
\pgfplotsset{compat=1.18}

\definecolor{desgreen}{HTML}{348443}
\DeclareMathOperator{\erfc}{\mathrm{erfc}}

\newcommand{\dx}{\,\mathrm{d}x}
\newcommand{\dt}{\,\mathrm{d}t}
\newcommand{\ddx}{\,\frac{\mathrm{d}}{\mathrm{d}x}}
\newcommand{\dydx}{\, \frac{\mathrm{d}y}{\mathrm{d}x}}
\newcommand{\du}{\, \mathrm{d}u}
\newcommand{\dv}{\, \mathrm{d}v}
\newcommand{\dw}{\, \mathrm{d}w}
\newcommand{\dy}{\, \mathrm{d}y}
\newcommand{\dr}{\, \mathrm{d}r}
\newcommand{\dtheta}{\, \mathrm{d}\theta}

\let\oldint\int
\renewcommand{\int}{\oldint\limits}

\delimiterfactor 1200




\newcounter{problem}
\newcommand{\problem}[1]{
	\refstepcounter{problem}
	\noindent\textbf{\theproblem}\quad #1
}

\newtoggle{solutions}
\togglefalse{solutions}

\newtoggle{studentpagebreaks}
\newtoggle{studentvspace}

\newtoggle{instructorpagebreaks}
\newtoggle{instructorvspace}

\iftoggle{solutions}{
	\togglefalse{studentpagebreaks}
	\togglefalse{studentvspace}
	\toggletrue{instructorpagebreaks}
	\toggletrue{instructorvspace}
}{
	\toggletrue{studentpagebreaks}
	\toggletrue{studentvspace}
	\togglefalse{instructorpagebreaks}
	\togglefalse{instructorvspace}
}

\newcommand{\spb}{
	\iftoggle{studentpagebreaks}{\newpage}{}
}
\newcommand{\svs}{
	\iftoggle{studentvspace}{\vspace*{0.5\textheight}}{}
}
\newcommand{\ipb}{
	\iftoggle{instructorpagebreaks}{\newpage}{}
}

\title{Example}
\author{By Kazuhiko Gushiken}
\date{\today}

\hypersetup{
	pdftitle={Example - Kazuhiko Gushiken},
	pdfauthor={Kazuhiko Gushiken},
	pdfsubject={Example Tex},
	pdfkeywords={example, latex}
}

\begin{document}
	\maketitle
	
	{\centering
		What is the Gamma Function? \\
		[12pt]
		The Gamma Function is defined as: 
		\begin{equation}
			\Gamma(z) = \int_{u=0}^{\infty} u^{z-1} e^{-u} \dt
		\end{equation}
		\text{Graph of }$\Gamma(z)$:\par
	}
	\vspace{-10pt}
	\begin{figure}[!hbtp]
		\centering
		\begin{tikzpicture}
			\begin{axis}[
				axis lines = center,
				xmin = -5,
				xmax = 5,
				ymin = -10,
				ymax = 10,
				minor tick num=4,
				minor tick style={draw=none},
				grid = both,
				width = \textwidth,
				height = 0.5\textwidth,
				axis line style={draw=black, thick},
				major tick length=5pt,
				ticklabel style={font=\bfseries},
				tick style={thick},
				xlabel = {$x$},
				ylabel = {$y$},
				unbounded coords=jump,
				clip=true,
				]
				\addplot[very thin, dashed, desgreen] coordinates {(0, 10) (0, -10)};
				\addplot[very thin, dashed, desgreen] coordinates {(-1, 10) (-1, -10)};
				\addplot[very thin, dashed, desgreen] coordinates {(-2, 10) (-2, -10)};
				\addplot[very thin, dashed, desgreen] coordinates {(-3, 10) (-3, -10)};
				\addplot[very thin, dashed, desgreen] coordinates {(-4, 10) (-4, -10)};
				\addplot[line width=1.2pt, red, samples=1600] gnuplot {
					(x < 0.2 && abs(sin(pi*x)) < 0.02) ? NaN : gamma(x)
					};
			\end{axis}
		\end{tikzpicture}
		\caption{Graph of the Gamma Function $\Gamma(z)$.}
	\end{figure}

\end{document}